\documentclass[11pt]{article}
\usepackage{fontspec,amsmath,amssymb,amsthm,hyperref}
\usepackage[margin=1in]{geometry}
\setlength{\emergencystretch}{3em}
\title{Spend one quarter of Fisher information on drift error}
\author{LangevinLean mathematical exposition}
\date{}
\begin{document}\maketitle
\noindent Source module: \texttt{UnadjustedLangevin/LSI/DriftError.lean}. Compile this proof document with XeLaTeX or LuaLaTeX. The accompanying HTML pairs the mathematical argument with the exact Lean statements and records the scope of the formalization.


Young’s inequality \(\langle a,b\rangle\le\frac14\|a\|^2+\|b\|^2\) bounds the error pairing in entropy production. Weighted energy integrability yields an actual integral inequality \(H\prime\le-\frac34J+\mathcal E\).


\section{inner le quarter norm sq add norm sq}


Cauchy–Schwarz gives \(\langle e,s\rangle\le\|e\|\|s\|\). Expand the nonnegative square \((\|e\|-\|s\|/2)^2\ge0\) to get \(\|e\|\|s\|\le\|e\|^2+\|s\|^2/4\). This is the Young coefficient that preserves three quarters of Fisher information.



\noindent Lean declarations: \texttt{inner\_le\_quarter\_norm\_sq\_add\_norm\_sq}.

\section{integrable weighted inner of energy}


Multiply the absolute Young bound by the nonnegative weight \(r\). It dominates \(|r\langle e,s\rangle|\) by \(r\|s\|^2/4+r\|e\|^2\). Both energy functions are integrable, so the assumed measurability and \texttt{Integrable.mono'} prove integrability of the cross term.



\noindent Lean declarations: \texttt{integrable\_weighted\_inner\_of\_energy}.

\section{entropy production le of forced weak identity}


First derive cross-term integrability from the two finite energies and its measurability. The forced weak identity tested by \(\log r+1\) gives production \(-I+\int r\langle e,s\rangle\). Insert the integrated Young bound to obtain \(-\frac34I+\int r\|e\|^2\). The entropy differential inequality is therefore proved from the stated spatial weak identity.



\noindent Lean declarations: \texttt{entropy\_production\_le\_of\_forced\_weak\_identity}.

\section{abs inner le quarter norm sq add norm sq}


Use the absolute Cauchy–Schwarz inequality \(|\langle e,s\rangle|\le\|e\|\|s\|\), followed by the same square expansion. The absolute-value version controls integrability of the signed cross term before its integral is estimated.



\noindent Lean declarations: \texttt{abs\_inner\_le\_quarter\_norm\_sq\_add\_norm\_sq}.

\section{integral weighted inner le quarter fisher add energy}


Multiply the pointwise Young inequality by \(r\ge0\), then integrate. The supplied integrability permits splitting the upper integral and pulling out \(1/4\). Thus \(\int r\langle e,s\rangle\le\frac14\int r\|s\|^2+\int r\|e\|^2\). No probability normalization of the weight is needed.



\noindent Lean declarations: \texttt{integral\_weighted\_inner\_le\_quarter\_fisher\_add\_energy}.

\end{document}
